% TOP_PROBLEM: {"id": 3137, "title": "3-Decomposition Conjecture", "category_id": 3, "category": {"id": 3, "name": "graph_theory", "display_name": "Graph Theory", "description": "Problems involving graphs, networks, and their properties.", "slug": "graph-theory", "order_index": 3, "created_at": "2026-07-31T15:26:25.671Z"}, "status": "open", "problem_number": "OPG-60027", "set_id": 12}
% FIRST_POSTED: 2026-10-01
% ATTEMPT_STATUS: unresolved
% UnsolvedMath ID 3137; source catalogue code OPG-60027
% !TeX program = lualatex
% Research attempt by GPT-6 Astra Ultra; no claim of a new complete solution.
\documentclass[11pt,a4paper]{article}
\usepackage[margin=25mm]{geometry}
\usepackage{fontspec}
\setmainfont{lmroman10-regular.otf}[BoldFont=lmroman10-bold.otf,
 ItalicFont=lmroman10-italic.otf,BoldItalicFont=lmroman10-bolditalic.otf]
\usepackage{amsmath,amssymb,mathtools}
\usepackage{unicode-math}
\setmathfont{latinmodern-math.otf}
\AtBeginDocument{\let\setminus\smallsetminus}
\usepackage{xurl,hyperref}
\hypersetup{colorlinks=true,urlcolor=blue}
\setcounter{secnumdepth}{0}
\setlength{\parindent}{0pt}
\setlength{\parskip}{5pt}
\setlength{\emergencystretch}{3em}
\begin{document}
\section*{3-Decomposition Conjecture}\label{3137}
{\small UnsolvedMath ID 3137 · OPG-60027 · Graph Theory · OpenGarden}
\subsection{Definitions and mathematical statement}
Let $G=(V,E)$ be a finite connected simple undirected graph in which
every vertex has degree three; such a graph is called cubic. A
spanning tree is a connected acyclic subgraph with vertex set $V$.
A matching is a set of edges no two of which share an endpoint. A
$2$-regular subgraph is a subgraph all of whose vertices have degree
two within that subgraph, so its components are vertex-disjoint
cycles.

The $3$-decomposition conjecture asks for an edge partition
\[
 E=E(T)\mathbin{\dot\cup}E(C)\mathbin{\dot\cup}M
\]
such that $T$ is a spanning tree of $G$, $C$ is a nonempty
$2$-regular subgraph of $G$, and $M$ is a matching, which is allowed
to be empty. Neither $C$ nor the vertices incident with $M$ need
span $G$. The dot on the unions requires that each edge belongs to
exactly one part. Vertices may occur in several of the subgraphs.

More explicitly, the conjecture quantifies over every connected
cubic $G$ and asks for a single choice of $T,C,M$ meeting all these
conditions. It does not merely ask for the three kinds of subgraph
to exist independently. Equivalently, one should choose a spanning
tree so that every nontrivial component of its edge complement is
either a cycle or a single edge. This complement description records
the matching condition on all remaining edges.
\subsection{Short English statement}
Can the edges of every connected cubic graph be partitioned into a spanning tree, disjoint cycles, and a matching?
\subsection{Research attempt}
\paragraph{Scope and process.}
This is a GPT-6 Astra Ultra research attempt dated 1 October 2026, using
primary-source browsing and elementary mathematical checks. The canonical
definition above is retained. Results below are restricted results or
reductions, not a claim of a new solution to the entire catalogue question.
No formal proof assistant or external mathematical referee checked this text.


\paragraph{Literature update.}
The historical conjecture is studied computationally in the primary
2025 paper [R1]. A preprint by Jicheng Ma posted in August 2026 [R2]
claims a full proof, including loopless cubic multigraphs. That claim
is reported here rather than silently treating the old catalogue
metadata as a current consensus. This attempt independently proves
an exact spanning-tree reformulation and an infinite family; it
does not certify every argument in the recent preprint.

\paragraph{The leaves of the spanning tree determine the cycle part.}
Suppose $(T,C,M)$ is a decomposition of a simple connected cubic
graph on $n\ge4$ vertices. At a vertex of $C$, two incident edges
are used by $C$. Since a spanning tree must give every vertex
positive degree, the third edge belongs to $T$, not $M$. Therefore
every vertex of $C$ is a leaf of $T$. Conversely a leaf of $T$
has two remaining incident edges. They cannot both lie in the
matching, and any vertex of a $2$-regular subgraph has either zero
or two edges of that subgraph. Hence both remaining edges belong
to $C$. Thus
\[
 V(C)=\{v:d_T(v)=1\}=:L.
\]
No edge of a tree on at least three vertices joins two leaves:
such an edge would form a component isolated from the other vertices.
Consequently every edge of $G$ with both endpoints in $L$ lies
outside $T$, and $G[L]=C$ is $2$-regular.

Conversely, suppose a spanning tree $T$ has leaf set $L$ such that
$G[L]$ is $2$-regular. Each leaf's two non-tree edges lie in
$G[L]$, exhausting those incidences. Each non-leaf has tree degree
two or three and therefore at most one non-tree edge; it has no
remaining edge to $L$. All non-tree edges outside $G[L]$ consequently
form a matching. This proves the exact equivalence:
\[
 \text{a 3-decomposition exists}\quad\Longleftrightarrow\quad
 \exists T\text{ spanning with }G[L(T)]\text{ 2-regular}.
\]
Such a leaf set is nonempty and has at least three vertices, so
the required nonempty cycle part is automatic.

\paragraph{Counting restrictions on a possible tree.}
Let $\ell$ be the number of tree leaves, and let $b$ be the number
of vertices of tree degree three. The tree degree sum gives
$\ell=b+2$. The remaining $n-\ell-b=n-2\ell+2$ vertices have tree
degree two, and each is an endpoint of exactly one matching edge.
Therefore
\[
 |M|=n/2-\ell+1,\qquad 3\le\ell\le n/2+1.
\]
The parity is consistent because every cubic graph has even order.
These identities are useful checks of a candidate decomposition;
they are not sufficient to produce the required adjacency pattern
among the leaves.

\paragraph{An explicit infinite family.}
For the prism $C_m\mathbin{\square}K_2$, $m\ge3$, name the two
cycles $u_0\cdots u_{m-1}u_0$ and $v_0\cdots v_{m-1}v_0$, and
join each $u_i$ to $v_i$. Put the whole $u$-cycle in $C$.
Put all $m$ vertical edges and the bottom path
$v_0v_1\cdots v_{m-1}$ in $T$. This graph is connected, has
$2m-1$ edges on $2m$ vertices, and is a path with one new leaf
attached at every vertex, hence is a tree. The single remaining
edge $v_{m-1}v_0$ forms $M$. Its leaf set is exactly the upper
cycle, verifying the reformulation directly. For $K_4$, take one
triangle as $C$ and the three edges to its fourth vertex as $T$;
the matching is empty, as allowed.

\paragraph{Why a Hamiltonian path does not solve the problem.}
Choosing a spanning path minimizes the number of leaves, but gives
exactly two. A simple $2$-regular graph cannot live on those two
vertices. Thus even in graphs with an obvious spanning path, that
choice can never be the tree part of the desired decomposition.
The task is to choose a different tree whose leaf adjacencies close
into cycles. Connectivity of a graph guarantees some spanning tree,
but supplies no such condition on its leaves.

\paragraph{Verification and final status.}
The script
\texttt{scripts/check\_attempt\_batch02\_algebra\_2026\_10\_01.py}
checks edge partitions, degrees and connectivity for the prism
construction for $3\le m\le20$, complementing the all-$m$ proof.
\textbf{UNRESOLVED as an independent complete proof}: the leaf
criterion is exact and the displayed family is settled. The
remaining independent step is construction of such a spanning tree
for every cubic graph, or verification of the complete claimed
argument in [R2].

\subsection{Sources}
\begin{itemize}
\item[S1] ULAM.AI and the UnsolvedMath Contributors, supplied problems.json, record 3137 (OPG-60027), OpenGarden collection. Catalogue identity and source transcription; CC BY 4.0.
\url{https://huggingface.co/datasets/ulamai/UnsolvedMath}.
\item[S2] Open Problem Garden, 3-Decomposition Conjecture. Original problem and discussion; the mathematical formulation below is expanded from the supplied source record.
\url{http://www.openproblemgarden.org/op/3_decomposition_conjecture}.

\item[R1] \emph{The 3-Decomposition Conjecture: A SAT-Based Approach
with Specialized Propagators}, CP 2025, LIPIcs 340, article 39.
Primary conference paper; scope checked 1 October 2026.
\url{https://drops.dagstuhl.de/entities/document/10.4230/LIPIcs.CP.2025.39}.
\item[R2] Jicheng Ma, \emph{Matching complements in subcubic graphs
and a proof of the 3-Decomposition Conjecture}, arXiv:2608.25385,
submitted 26 August 2026. Reported proof claim; not independently
certified in this attempt.
\url{https://arxiv.org/abs/2608.25385}.

\end{itemize}
\end{document}
